\documentclass{article}

% Language setting
% Replace `english' with e.g. `spanish' to change the document language
\usepackage[english]{babel}
\usepackage{float}
\usepackage{array}
% Set page size and margins
% Replace `letterpaper' with `a4paper' for UK/EU standard size
\usepackage[letterpaper,top=1cm,bottom=2cm,left=2cm,right=2cm,marginparwidth=1.75cm]{geometry}

\usepackage{longtable} 
\usepackage{multicol}
\usepackage[dvipsnames]{xcolor}
% Useful packages 
\usepackage{bbding}
\usepackage{pifont}
\usepackage{wasysym}

\usepackage{mathtools}
\usepackage{multirow}
\usepackage{amsmath}
\usepackage{graphicx}
\usepackage[colorlinks=true, allcolors=blue]{hyperref}
\usepackage{makeidx}
\setlength{\columnsep}{1cm}
\usepackage{xcolor}
\usepackage{graphicx}
\usepackage{tcolorbox}
\usepackage{algpseudocode} 
\usepackage{algorithm}
\algnewcommand\algorithmicnot{\textbf{not}}
\algdef{SE}[IF]{IfNot}{EndIf}[1]{\algorithmicif\ \algorithmicnot\ #1\ \algorithmicthen}{\algorithmicend\ \algorithmicif}%

\tcbuselibrary{skins}
\usepackage{lipsum}
\definecolor{blueboxTitle}{HTML}{B3C5D7}
\definecolor{blueboxFrame}{HTML}{7392B7}
\definecolor{blueboxBackground}{HTML}{C5D5EA}
\definecolor{rose}{HTML}{CC7E85}
\definecolor{rose-d}{HTML}{83343C}
\definecolor{rose-l}{HTML}{DCA7AC}
\definecolor{boxBackground}{HTML}{DEF4C6}
\definecolor{boxFrame}{HTML}{004346}
\definecolor{boxTitle}{HTML}{B1CF5F}
\definecolor{boxBackground1}{HTML}{F0D3F7}
\definecolor{boxFrame1}{HTML}{A57982}
\definecolor{boxTitle1}{HTML}{B98EA7}
\tcbset{my box1/.style={
    enhanced, fonttitle=\bfseries,
    colback=white, colframe=boxFrame1,
    coltitle=black, colbacktitle=boxTitle1
  }
}
\tcbset{my info/.style={
    enhanced, fonttitle=\bfseries,
    colback=white, colframe=blueboxFrame,
    coltitle=black, colbacktitle=blueboxTitle
  }
}
\tcbset{my box/.style={
    enhanced, fonttitle=\bfseries,
    colback=white, colframe=boxFrame,
    coltitle=black, colbacktitle=boxTitle
  } 
}
\tcbset{my other box/.style={
    enhanced, fonttitle=\bfseries,
    colback=white, colframe=rose-d,
    coltitle=black, colbacktitle=rose
  }
}
\newtcolorbox{definitionbox}[1][]{my info, #1}
\newtcolorbox{action}[1][]{my box, #1}
\newtcolorbox{vsbox}[1][]{my box1, #1}
\newtcolorbox{algo}[1][]{my other box, #1}
\newcommand{\checkedbox}[0]{\makebox[0pt][l]{$\square$}\raisebox{.15ex}{\hspace{0.1em}$\checkmark$}}

\usepackage{amssymb} 
\usepackage{biblatex} %Imports biblatex package
\addbibresource{sample.bib} %Import the bibliography file

\begin{document}
\begin{align*}
    \text{maximiere} \quad & z = 2x_1 - 3x_2 \\
    \text{so dass} \quad & x_3 = 2 - x_1 + x_2 \\
                          & x_4 = 3 + 2x_1 - 2x_2 \\
                          & x_1, x_2, x_3, x_4 \geq 0
\end{align*}

Die Eingangsvariable könnten $x_1$ oder $x_2$ sein, da: 
\[ 
z = 2\textcolor{Peach}{x_1} + 3\textcolor{Peach}{x_2 }
\]
Wir wählen $x_1$ als Entering-Variable. \newline
Jetzt müssen wir berechnen was der Wert (die Ratio) von dem \textcolor{RedViolet}{Koeffizienten} der Entering-Variable durch das jeweilige "\textcolor{BlueViolet}{b}" der Gleichung ist. 
\begin{align*}
 x_3 &= \textcolor{BlueViolet}{2} - \textcolor{RedViolet}{1}x_1 + x_2 \\
x_4 & = \textcolor{BlueViolet}{3} + \textcolor{RedViolet}{2}x_1 - 2x_2 \\
\end{align*}
Die Ratios sind: 
\begin{align*}
    \frac{2}{1} &= 0.5 & \quad \text{für } x_3 \\
    \frac{3}{2} & \approx 0.6 & \quad  \text{ für } x_3 \\
\end{align*}
Die Leaving-Variable wird die mit der kleinsten Ratio, also $x_3$. \newline
Nun formen wir die zur Leave-Variable ($x_3$) gehörende Gleichung nach der Entering-Variable ($x_1$) um.
\begin{align*}
    x_3 &= 2 - x_1 + x_2 \\
    x_1 + x_3 &= 2 + x_2 \\
    x_1  &= 2 + x_2  - x_3 \\
\end{align*}
Nun müssen wir noch das LP updaten mit der neuen Variable:
\begin{align*}
    z &= 2x_1 - 3x_2 \\
    &= 2(2 + x_2  - x_3) - 3x_2 \\
    &= 4 + 2x_2  - 2x_3 - 3x_2 \\
    &= 4 + 2x_2  - 5x_2 \\
\end{align*}
\begin{align*}
     x_4 &= 3 + 2(2 + x_2  - x_3) - 2x_2 \\
     &= 3 + 4 + 2x_2  - 2x_3 - 2x_2 \\
     &= 7 - 2x_3 \\
\end{align*}
Das LP sieht nach dem Pivotschritt also folgendermaßen aus:
\begin{align*}
    \text{maximiere} \quad & z = 4 + 2x_2  - 5x_2 \\
    \text{so dass} \quad & x_1  = 2 + x_2  - x_3 \\
                          & x_4 = 7 - 2x_3 \\
                          & x_1, x_2, x_3, x_4 \geq 0
\end{align*}

\end{document}